\documentclass[10pt,a4paper]{article}

\usepackage[T1]{fontenc}
\usepackage[utf8]{inputenc}
\usepackage[ngerman]{babel}
\usepackage[a4paper,margin=14mm]{geometry}
\usepackage{lmodern}
\usepackage[table]{xcolor}
\usepackage{tabularx}
\usepackage{array}
\usepackage{microtype}
\usepackage[most]{tcolorbox}
\usepackage{tikz}
\usetikzlibrary{arrows.meta,calc,positioning}

\pagestyle{empty}
\setlength{\parindent}{0pt}
\setlength{\parskip}{0pt}
\renewcommand{\familydefault}{\sfdefault}
\renewcommand{\arraystretch}{1.18}

\definecolor{Navy}{HTML}{17324D}
\definecolor{Blue}{HTML}{3B6EA5}
\definecolor{Sky}{HTML}{EAF3FA}
\definecolor{Mint}{HTML}{E8F5EF}
\definecolor{Gold}{HTML}{F3B33D}
\definecolor{Ink}{HTML}{1E2935}
\definecolor{Muted}{HTML}{617080}
\definecolor{Line}{HTML}{D8E1E8}
\definecolor{CodeBg}{HTML}{F5F7F9}
\definecolor{KeyBg}{HTML}{FFF3D6}
\definecolor{KeyLine}{HTML}{B87810}

\color{Ink}

\newcommand{\sectiontitle}[1]{%
  \vspace{2.2mm}
  {\color{Blue}\bfseries\large #1}\par
  \vspace{0.7mm}
  {\color{Blue}\rule{\linewidth}{0.7pt}}\par
  \vspace{1.1mm}
}

\newcommand{\dbtable}[4]{%
  \begin{tabular}{@{}#1@{}}
    \rowcolor{Sky}#2\\
    #3\\
    #4
  \end{tabular}%
}

\begin{document}

\colorbox{Navy}{%
  \parbox{\dimexpr\linewidth-2\fboxsep\relax}{%
    \vspace{3mm}\color{white}
    \begin{tabularx}{\linewidth}{@{}Xr@{}}
      {\small\bfseries INFORMATIK · KLASSE 10} &
      \colorbox{Gold}{\color{Navy}\small\bfseries\strut LÖSUNGEN}
    \end{tabularx}
    \vspace{2.4mm}

    {\fontsize{21}{24}\selectfont\bfseries 1:n-Beziehungen und Fremdschlüssel}\par
    \vspace{0.8mm}
    {\large Beispiel: InstaHub -- \texttt{users} und \texttt{photos}}
    \vspace{2.8mm}
  }%
}

\vspace{2.5mm}
\colorbox{Sky}{%
  \parbox{\dimexpr\linewidth-2\fboxsep\relax}{%
    \vspace{1.6mm}
    \textcolor{Blue}{\bfseries Ziel:}
    Du erkennst eine 1:n-Beziehung und erklärst, wie ein Fremdschlüssel zwei
    Tabellen über passende IDs verbindet.
    \vspace{1.6mm}
  }%
}

\sectiontitle{1 · Ein Benutzer, mehrere Fotos}

\begin{center}
\begin{tabular}{@{}c@{\hspace{6mm}}c@{}}
  \begin{minipage}[t]{48mm}
  \centering
  \textcolor{Blue}{\bfseries\texttt{users}}\par\vspace{0.8mm}
  \begin{tabular}{|>{\ttfamily\small}r|>{\ttfamily\small}l|}\hline
    \rowcolor{Sky}\textbf{id} & \textbf{username}\\\hline
    \cellcolor{KeyBg}\textbf{3} & bergcoder\\\hline
  \end{tabular}
  \end{minipage}
  &
  \begin{minipage}[t]{96mm}
  \centering
  \textcolor{Blue}{\bfseries\texttt{photos}}\par\vspace{0.8mm}
  \begin{tabular}{|>{\ttfamily\small}r|>{\ttfamily\small}r|>{\ttfamily\small}p{37mm}|}\hline
    \rowcolor{Sky}\textbf{id} & \textbf{user\_id} & \textbf{description}\\\hline
    11 & \cellcolor{KeyBg}\textbf{3} & Heute draußen im Wald programmiert …\\\hline
    12 & \cellcolor{KeyBg}\textbf{3} & Bin gerade am Debuggen mitten im Wald …\\\hline
    13 & \cellcolor{KeyBg}\textbf{3} & Ein neuer Pfad, ein neuer Gedanke …\\\hline
  \end{tabular}
  \end{minipage}
\end{tabular}
\end{center}

\vspace{1.7mm}
\begin{center}
\begin{tikzpicture}[
  classcard/.style={draw=Blue, line width=0.8pt, minimum width=49mm, minimum height=20mm, inner sep=0pt}
]
  \node[classcard] (usersclass) {};
  \node[classcard, right=58mm of usersclass] (photosclass) {};
  \draw[Blue, line width=0.8pt]
    ([yshift=-7mm]usersclass.north west) -- ([yshift=-7mm]usersclass.north east);
  \draw[Blue, line width=0.8pt]
    ([yshift=-7mm]photosclass.north west) -- ([yshift=-7mm]photosclass.north east);
  \draw[Blue, line width=0.8pt]
    ([yshift=-7mm]usersclass.north east) -- ([yshift=-7mm]photosclass.north west);
  \node[font=\large\bfseries\ttfamily] at ([yshift=-3.5mm]usersclass.north) {users};
  \node[font=\large\bfseries\ttfamily] at ([yshift=-3.5mm]photosclass.north) {photos};
  \node[font=\large\bfseries] at ([xshift=9mm,yshift=-4mm]usersclass.north east) {1};
  \node[font=\large\bfseries] at ([xshift=-9mm,yshift=-4mm]photosclass.north west) {n};
\end{tikzpicture}\\[-0.3mm]
{\small\color{Muted}Ein Benutzer kann mehreren Fotos zugeordnet sein.}
\end{center}

\begin{tcolorbox}[colback=white,colframe=Line,arc=2mm,boxrule=0.7pt,left=2.5mm,right=2.5mm,top=1.5mm,bottom=1.5mm]
  \textbf{1:n-Beziehung:} Bei einer 1:n-Beziehung kann einem Datensatz der
  einen Tabelle mehr als ein Datensatz der anderen Tabelle zugeordnet sein.
  \par\smallskip
  Beispiel: \textbf{Ein Benutzer kann mehrere Fotos hochladen. Jedes Foto gehört
  zu genau einem Benutzer.}
\end{tcolorbox}

\sectiontitle{2 · Der Fremdschlüssel stellt die Verbindung her}

\begin{center}
\begin{tikzpicture}[>=Latex, every node/.style={font=\small},
  key/.style={draw=KeyLine, fill=KeyBg, rounded corners=1.2pt, inner sep=2.4pt, line width=0.8pt},
  note/.style={draw=Blue, fill=Sky, rounded corners=1.2pt, inner sep=2pt, align=center}]
  \node[key] (uid) {\texttt{users.id = 3}};
  \node[key, right=42mm of uid] (fk) {\texttt{photos.user\_id = 3}};
  \draw[->, line width=0.9pt, Blue] (fk) -- node[above, fill=white, inner sep=1.3pt] {verweist auf} (uid);
  \node[note, below=7mm of $(uid)!0.5!(fk)$] (meaning) {\textbf{Verbindung über die Benutzer-ID}\\[-1pt]{\scriptsize Das Foto gehört zu dem Benutzer mit id 3.}};
\end{tikzpicture}
\end{center}

\vspace{0.4mm}
\begin{tabularx}{\linewidth}{@{}p{38mm}X@{}}
  \colorbox{KeyBg}{\parbox{34mm}{\centering\textbf{Fremdschlüssel}\par\small\texttt{photos.user\_id}}}
  &
  \textbf{Ein Fremdschlüssel ist ein Attribut, das auf den Primärschlüssel
  eines Datensatzes in einer anderen Tabelle verweist und dadurch die Tabellen
  miteinander verbindet.}
\end{tabularx}

\vspace{1.4mm}
\begin{tabularx}{\linewidth}{@{}X X@{}}
  \textcolor{Blue}{\bfseries\texttt{users.id}}\par
  identifiziert den Benutzer und ist hier der Primärschlüssel. &
  \textcolor{Blue}{\bfseries\texttt{photos.user\_id}}\par
  enthält dessen ID und ist der Fremdschlüssel.
\end{tabularx}

\vspace{1.7mm}
\colorbox{Mint}{%
  \parbox{\dimexpr\linewidth-2\fboxsep\relax}{%
    \vspace{1.5mm}
    \textcolor{Blue}{\bfseries Merke:}
    Eine Benutzer-ID kann in der Tabelle \texttt{photos} in mehreren Zeilen
    vorkommen. Jedes einzelne Foto verweist mit genau einer
    \texttt{user\_id} auf seinen Benutzer.
    \vspace{1.5mm}
  }%
}

\vfill
{\color{Line}\rule{\linewidth}{0.5pt}}\par
\vspace{1.3mm}
\begin{tabularx}{\linewidth}{@{}Xr@{}}
  {\small\color{Muted}Klasse 10 · Datenbanken} &
  {\small\color{Muted}Sicherungsblatt · Aufgabe 3}
\end{tabularx}

\end{document}
